\vfill%
\pagebreak

\section*{Exercises}
\markright{\MakeUppercase{Exercises}}%
\subsection*{Exercise 1}

Write a program that prints the numbers from 1 to 15, one per line, except that
multiples of 3 are replaced by \textit{Fizz}, multiples of 5 by \textit{Buzz},
and multiples of both by \textit{FizzBuzz}.

\subsection*{Exercise 2}

What does the following program print? Work it out by hand, iteration by
iteration, then check your answer by compiling and running it.

\begin{lstlisting}[style=coloredverbatim]
use std::io;

fn main() {
  let mut i = 0;
  while i < 10 {
    i += 1;
    if i % 3 == 0 {
      continue;
    }
    if i == 8 {
      break;
    }
    print(i, " ");
  }
  println();
  println("i = ", i);
}
\end{lstlisting}

\subsection*{Exercise 3}

The following program finds the smallest number whose square is at least the
number typed by the user, and says whether the number typed is negative. It contains
three errors. Find them, correct them, and compile the program to test your
corrections.

\begin{lstlisting}[style=coloredverbatimError, caption=Ymir program with errors]
use std::io;

fn main() {
  let n = read!i32(ask-> "A number: ");
  let mut root = 0;
  while true {
    if root * root >= n {
      break;
    }
    root += 1;
  }

  let sign = if n < 0 { "negative" };
  if root {
    println("Smallest root: ", root);
  }
  println(n, " is ", sign);
}
\end{lstlisting}

\subsection*{Exercise 4}

Write a program that reads a positive number and prints the sum of its digits.
For \token{40953}, it prints \token{21}. (\textit{Hint: the remainder of a
division by ten is the last digit of a number, and the division by ten removes
that digit.})

\subsection*{Exercise 5}

Write a program that reads a height and draws a triangle of stars that many rows
high. For example, with a height of 4:

\begin{lstlisting}[style=bashVerb]
*
**
***
****
\end{lstlisting}

\subsection*{Exercise 6}

Write a program that reads a month number and prints how many days that month
has, ignoring leap years, or a message if there is no such month. Compute the
number of days with a single \token{match} used as a value. In what order must
the arms be written, and why?

\subsection*{Exercise 7}

Rewrite the second loop of Listing~\ref{lst:control:for_reverse}, which prints
\token{0 5 10 15 20}, with a \token{while} loop instead of a \token{for}.

\subsection*{Exercise 8}

Write a program that reads a number and tells whether it is prime, that is,
whether it is at least 2 and has no divisor other than 1 and itself.

\subsection*{Exercise 9}

Modify the guessing game of Listing~\ref{lst:control:guessing_game} so that the
player has only seven tries. After seven wrong guesses, the program stops and
reveals the number. (\textit{Hint: make the \token{loop} produce a \token{bool}
that says whether the number was found.})

\vfill%
\pagebreak

\forceevenpage{}

\section*{Solutions}
\markright{\MakeUppercase{Solutions}}%
\subsection*{Exercise 1}

The test for \textit{FizzBuzz} must come first: a multiple of 15 is also a
multiple of 3, and the chain stops at the first condition that holds.

\begin{lstlisting}[style=coloredverbatim, caption=Solution for exercise 1]
use std::io;

fn main() {
  for i in 1 ... 15 {
    if i % 15 == 0 {
      println("FizzBuzz");
    } else if i % 3 == 0 {
      println("Fizz");
    } else if i % 5 == 0 {
      println("Buzz");
    } else {
      println(i);
    }
  }
}
\end{lstlisting}

\subsection*{Exercise 2}

\begin{lstlisting}[style=bashVerb]
1 2 4 5 7
i = 8
\end{lstlisting}

3 and 6 are skipped by \token{continue}, which returns to the test of the
\token{while} before \token{print} is reached. When \token{i} reaches 8,
\token{break} leaves the loop, also before \token{print}, and \token{i} keeps the
value 8.

\subsection*{Exercise 3}

\token{while true} is not allowed, and is written \token{loop}. The \token{if}
that gives \token{sign} its value has no \token{else}, so there is no value when
\token{n} is not negative. And \token{root} is an integer, not a Boolean, so the
condition of the last \token{if} must be a comparison.

\begin{lstlisting}[style=coloredverbatim, escapechar=@, caption=Solution for exercise 3]
use std::io;

fn main() {
  let n = read!i32(ask-> "A number: ");
  let mut root = 0;
  @\hcb{loop}@ {
    if root * root >= n {
      break;
    }
    root += 1;
  }

  let sign = if n < 0 { "negative" } @\hcb{else \{ "positive or zero" \}}@;
  if @\hcb{root != 0}@ {
    println("Smallest root: ", root);
  }
  println(n, " is ", sign);
}
\end{lstlisting}

\subsection*{Exercise 4}

\begin{lstlisting}[style=coloredverbatim, caption=Solution for exercise 4]
use std::io;

fn main() {
  let mut n = read!i32(ask-> "A positive number: ");
  let mut sum = 0;
  while n > 0 {
    sum += n % 10;
    n /= 10;
  }
  println("Sum of the digits: ", sum);
}
\end{lstlisting}

\subsection*{Exercise 5}

The outer loop picks the row, the inner loop prints one star per column. The
inner iterator is never read, so it is \token{\_}.

\begin{lstlisting}[style=coloredverbatim, caption=Solution for exercise 5]
use std::io;

fn main() {
  let height = read!i32(ask-> "Height: ");
  for row in 1 ... height {
    for _ in 0 .. row {
      print("*");
    }
    println();
  }
}
\end{lstlisting}

\subsection*{Exercise 6}

The arms are tried from top to bottom. The arm \token{1 ... 12} also fits
February and the months of 30 days, so it must come after them, and the
catch-all \token{\_}, which fits everything, must come last. It is also required:
without it, the \token{match} would have no value for 13.

\begin{lstlisting}[style=coloredverbatim, caption=Solution for exercise 6]
use std::io;

fn main() {
  let month = read!i32(ask-> "Month (1 to 12): ");
  let days = match month {
    2 => 28
    4 | 6 | 9 | 11 => 30
    1 ... 12 => 31
    _ => 0
  };

  if days == 0 {
    println("There is no month ", month, ".");
  } else {
    println("Month ", month, " has ", days, " days.");
  }
}
\end{lstlisting}

\subsection*{Exercise 7}

The iterator becomes a mutable variable declared before the loop, the bound
becomes the condition, and the step becomes an addition at the end of the block.
The \token{for} did all three by itself, which is why it is the better choice
when the numbers are known in advance.

\begin{lstlisting}[style=coloredverbatim, caption=Solution for exercise 7]
use std::io;

fn main() {
  let mut i = 0;
  while i <= 20 {
    print(i, " ");
    i += 5;
  }
  println();
}
\end{lstlisting}

\subsection*{Exercise 8}

The program assumes the number is prime, then looks for a divisor between 2 and
the number itself, excluded. The first divisor found proves the assumption
wrong, and there is no need to look further. A number below 2 is not prime, and
the range \token{2 .. n} is then empty, so the loop does not run at all.

\begin{lstlisting}[style=coloredverbatim, caption=Solution for exercise 8]
use std::io;

fn main() {
  let n = read!i32(ask-> "A number: ");
  let mut prime = n >= 2;
  for d in 2 .. n {
    if n % d == 0 {
      prime = false;
      break;
    }
  }

  if prime {
    println(n, " is prime.");
  } else {
    println(n, " is not prime.");
  }
}
\end{lstlisting}

The loop tests more divisors than it needs to. A divisor of \token{n} other
than \token{n} itself is at most half of \token{n}, so the range can stop there,
\token{2 ... n / 2}, and the program does half the work for the same answer.

\subsection*{Exercise 9}

The loop has two ways out, and each \token{break} says which one was taken.
Checking the number of tries at the top of the block, before reading, stops the
game after the seventh wrong guess without asking for an eighth.

\begin{lstlisting}[style=coloredverbatim, caption=Solution for exercise 9]
use std::io;
use std::rand;

fn main() {
  let secret = uniform(1, 100);
  let mut tries = 0;

  println("I am thinking of a number between 1 and 100. You have 7 tries.");
  let found = loop {
    if tries == 7 {
      break false;
    }

    let guess = read!i32(ask-> "Your guess: ");
    tries += 1;

    if guess < secret {
      println("Too small.");
    } else if guess > secret {
      println("Too big.");
    } else {
      break true;
    }
  };

  if found {
    println("Found it in ", tries, " tries!");
  } else {
    println("Lost! The number was ", secret, ".");
  }
}
\end{lstlisting}
